\uuid{O2jD}
\exo7id{7460}
\titre{exo7 7460}
\auteur{mourougane}
\organisation{exo7}
\datecreate{2021-08-10}
\isIndication{false}
\isCorrection{false}
\chapitre{Géométrie affine dans le plan et dans l'espace}
\sousChapitre{Géométrie affine dans le plan et dans l'espace}
\cpge{mpsi-mp2i.applications-lineaires.b, pcsi.applications-lineaires.b, ptsi.applications-lineaires.b, mpsi-mp2i.sous-espaces-affines, ptsi.espaces-vectoriels.b, mp-mpi.endomorphismes-euclidien.c, pc.endomorphismes-euclidiens.a, psi.endomorphismes-euclidien.a, pt.isometries-euclidien.a}
\module{Géométrie}
\niveau{L2}
\difficulte{}

\contenu{
\question{
On munit le plan affine euclidien d'un repère orthonormé $(O,i,j)$ et on l'identifie au plan complexe.
\'Ecrire à l'aide des affixes complexes, la symétrie glissée d'axe d'équation $x+y=2$ et de vecteur $(3,3)$.
}
}
